\subsection{Native-band benefits survive support and asymmetry controls}\label{sec:pvnew}
The PV controls test whether impossible negative outputs or symmetric intervals explain
the earlier benefit. After both constructions are separately tuned using their projected
validation scores, CQR's mean gain over symmetric residual intervals falls from 20.51\%
to 15.61\% overall. Against the projected signed-residual comparator it is 10.86\%
(Table~\ref{tab:pvcontrols}). The corresponding positive-output gains are 12.77\%,
11.57\% and 9.56\%. The stronger controls reduce but do not remove the native-band
benefit. Because the denominators and selected settings differ, successive percentage
changes are not additive causal contributions.

\begin{table}[!htbp]\centering\scriptsize
\caption{PV construction controls, 32 pairs, validation-tuned settings. Gains are
paired CQR improvements over the named baseline (\%); coverages are percent. High
forecast denotes the fixed prefix-median point-forecast split, known at issue time.}
\label{tab:pvcontrols}\begin{tabular}{lllrrr}
\toprule
Support & Baseline & Group & CQR gain & Cov. CQR & Cov. base \\
\midrule
Raw & Residual & all & 20.51 & 89.72 & 89.34 \\
Raw & Residual & positive & 12.77 & 86.32 & 80.17 \\
Raw & Residual & forecast high & 14.70 & 87.03 & 82.91 \\
Nonnegative & Residual & all & 15.61 & 91.56 & 90.55 \\
Nonnegative & Residual & positive & 11.57 & 86.45 & 80.50 \\
Nonnegative & Residual & forecast high & 12.41 & 87.07 & 83.04 \\
Nonnegative & Signed & all & 10.86 & 91.56 & 91.40 \\
Nonnegative & Signed & positive & 9.56 & 86.45 & 80.09 \\
Nonnegative & Signed & forecast high & 9.26 & 87.07 & 84.59 \\
\bottomrule
\end{tabular}

\end{table}

Projection raises selected CQR overall coverage to 91.56\%, partly by including
previously missed zero outputs. It does not solve positive-output coverage (86.45\%).
The high-forecast group, averaging 53.22\% of test observations, has a 9.26\% score gain
against projected signed residuals but only 87.07\% CQR coverage. Undercoverage is
therefore also visible in an issue-time observable subgroup, without routing by the
future target. Width and tail-penalty components, raw/projected defaults, all candidate
settings and the separately selected rolling controls are supplied for every pair.
